% !PRINTED 76
\weberpage{76}

2. Da zwei der $\vartheta$-Quadrate (12) nicht in konstantem Verhältnis stehen, so besteht auch keine Gleichung von der Form $F^{(n)}(\Theta_0,\Theta_1) = 0$.

3. Die Gesamtzahl der Konstanten, die nach (16), (17) in den zu einer und derselben Charakteristik gehörenden geraden und ungeraden Funktionen auftreten, ist genau gleich der Ordnung $m$.

\begin{center}
\textbf{\S\,22. \gesperrt{Das Additionstheorem.}}
\end{center}

Sind $u$, $v$ zwei Veränderliche, so gehören die Produkte
\[
\vartheta_{g_1,g_2}(u+v)\,\vartheta_{g'_1,g'_2}(u-v)
\]
zu den $\Theta$-Funktionen zweiter Ordnung, mit der Charakteristik
\[
(g_1+g'_1,\,g_2+g'_2),
\]
und zwar für jede der beiden Variablen $u$, $v$; sie sind also als Funktionen von $v$ linear darstellbar durch die Funktionen (12) und (15) in \S\,21, so daß die Variable $u$ in den Koeffizienten vorkommt. Diese Darstellung ist, wenn die Charakteristik $(0,0)$ ist, auf mehrfache Art möglich, da man zwei beliebige der Funktionen (12) wählen kann, in den anderen Fällen nur auf eine Art. Man erhält so 16 Formeln, von denen aber 6 durch bloße Vertauschung von $v$ mit $-v$ aus den anderen herzuleiten sind, so daß nur 10 wesentlich verschiedene bleiben. Man leitet diese Formeln sehr leicht ab, indem man sie zunächst mit unbestimmten Koeffizienten ansetzt und diese dann dadurch bestimmt, daß man für $v$ solche spezielle Werte setzt, für die je eine der Funktionen $\vartheta(v)$ verschwindet. So ist z.\,B.:
\[
\vartheta_{00}(u+v)\,\vartheta_{00}(u-v) = A\,\vartheta_{00}^2(v) + B\,\vartheta_{11}^2(v),
\]
und wenn man $v=0$ und $v=\dfrac{1+\omega}{2}$ setzt, so erhält man mit Benutzung der Formeln des \S\,21:
\[
A\,\vartheta_{00}^2 = \vartheta_{00}^2(u), \quad B\,\vartheta_{00}^2 = \vartheta_{11}^2(u),
\]
also

\weq{1}{\vartheta_{00}^2\,\vartheta_{00}(u+v)\,\vartheta_{00}(u-v) = \vartheta_{00}^2(u)\,\vartheta_{00}^2(v) + \vartheta_{11}^2(u)\,\vartheta_{11}^2(v),}
und wenn man hierin $u$ durch
\[
u+\frac{1}{2}, \quad u+\frac{\omega}{2}, \quad u+\frac{1+\omega}{2}
\]
ersetzt, so bildet man drei weitere Relationen, die man auf demselben Wege wie (1) auch direkt hätte ableiten können:
%CONT>>
